\chapter{Hypothesis Testing}

\section{Basic Setup}

Let $\theta \in \Omega$ be a parameter vector and define restrictions
\[
R(\theta) = [R_1(\theta), \dots, R_J(\theta)]'
\]

\[
\Omega_0 = \{\theta \in \Omega : R(\theta)=0\}, \quad
\Omega_1 = \{\theta \in \Omega : R(\theta)\neq 0\}
\]

Hypotheses:
\[
H_0: \theta \in \Omega_0
\quad
H_1: \theta \in \Omega_1
\]

\section{Errors}

\begin{itemize}
    \item Type I error:
    \[
    P(\text{reject } H_0 \mid H_0 \text{ true}) = \alpha
    \]
    \item Type II error:
    \[
    P(\text{fail to reject } H_0 \mid H_1 \text{ true})
    \]
\end{itemize}

\section{Likelihood Ratio Test}

\[
\lambda = \frac{l^*(\Omega)}{l^*(\Omega_0)}
\]

Decision rule:
\[
\text{Reject } H_0 \text{ if } \lambda \ge \lambda_0
\]
$\lambda_0$ is a constant chosen on the basis of the relative cost of type-1 versus type-2 errors.

\subsection{p-value Rule}

\[
p = P(\lambda \ge \lambda_T \mid H_0)
\]

\[
\text{Reject } H_0 
\]

\[
\text{ if } \lambda_T \ge \lambda_0
\]

\[
hence, \quad p \le \alpha
\]


\section{Example 1: Mean Test}

Let $(Y_1,\dots,Y_T)$ be a random sample from
\[
Y_t \sim N(\beta,\sigma^2).
\]

Consider the hypotheses
\[
H_0:\beta=\beta_0,
\qquad
H_1:\beta\neq\beta_0.
\]

\paragraph{Likelihood function.}

The likelihood function of the sample is
\[
L(\beta,\sigma^2)
=
(2\pi\sigma^2)^{-T/2}
\exp\!\left(
-\frac{1}{2\sigma^2}
\sum_{t=1}^T (Y_t-\beta)^2
\right).
\]

The maximum likelihood estimators are
\[
\hat{\beta} = \frac{1}{T}\sum_{t=1}^T Y_t,
\qquad
\hat{\sigma}^2 = \frac{1}{T}\sum_{t=1}^T (Y_t-\hat{\beta})^2.
\]

Under $H_0$, define
\[
\hat{\sigma}_0^2 = \frac{1}{T}\sum_{t=1}^T (Y_t-\beta_0)^2.
\]

The maximized likelihoods are
\[
l^*(\Omega)
=
(2\pi\hat{\sigma}^2)^{-T/2} e^{-T/2},
\qquad
l^*(\Omega_0)
=
(2\pi\hat{\sigma}_0^2)^{-T/2} e^{-T/2}.
\]


\paragraph{Simplification of the likelihood.}

Note that
\[
\hat{\sigma}^2
=
\frac{1}{T}\sum_{t=1}^T (Y_t-\hat{\beta})^2.
\]

Hence,
\[
\sum_{t=1}^T \frac{(Y_t-\hat{\beta})^2}{\hat{\sigma}^2}
=
\frac{\sum_{t=1}^T (Y_t-\hat{\beta})^2}
{\frac{1}{T}\sum_{t=1}^T (Y_t-\hat{\beta})^2}
=
T.
\]



\paragraph{Likelihood ratio.}

The likelihood ratio is
\[
\lambda
=
\frac{l^*(\Omega)}{l^*(\Omega_0)}
=
\left(
\frac{\hat{\sigma}_0^2}{\hat{\sigma}^2}
\right)^{T/2}.
\]

Since the distribution of $\lambda$ is not convenient, consider the monotonic transformation
\[
\lambda^*
=
(T-1)\left(\lambda^{2/T} - 1\right).
\]
After simplification, we obtain
\begin{align*}
\lambda^*
&=
(T-1)\frac{\hat{\sigma}_0^2-\hat{\sigma}^2}{\hat{\sigma}^2} \\
&=
(T-1)
\frac{
\frac{1}{T}\sum_{t=1}^T (Y_t-\beta_0)^2
-
\frac{1}{T}\sum_{t=1}^T (Y_t-\hat{\beta})^2
}{
\frac{1}{T}\sum_{t=1}^T (Y_t-\hat{\beta})^2
} \\
&=
(T-1)
\frac{
\sum_{t=1}^T (Y_t-\beta_0)^2
-
\sum_{t=1}^T (Y_t-\hat{\beta})^2
}{
\sum_{t=1}^T (Y_t-\hat{\beta})^2
}.
\end{align*}

Expanding the squares in the numerator,
\begin{align*}
\sum_{t=1}^T 
&(Y_t-\beta_0)^2
-
\sum_{t=1}^T (Y_t-\hat{\beta})^2\\
&=
\sum_{t=1}^T
\Bigl[
(Y_t^2 - 2Y_t\beta_0 + \beta_0^2)
-
(Y_t^2 - 2Y_t\hat{\beta} + \hat{\beta}^2)
\Bigr] \\
&=
\sum_{t=1}^T
\Bigl[
-2Y_t\beta_0 + \beta_0^2 + 2Y_t\hat{\beta} - \hat{\beta}^2
\Bigr].
\end{align*}

Using
\[
\hat{\beta}=\frac{1}{T}\sum_{t=1}^T Y_t,
\qquad\text{so that}\qquad
\sum_{t=1}^T Y_t = T\hat{\beta},
\]
we get
\begin{align*}
\sum_{t=1}^T 
&(Y_t-\beta_0)^2
-
\sum_{t=1}^T (Y_t-\hat{\beta})^2\\
&=
-2\beta_0 \sum_{t=1}^T Y_t
+ T\beta_0^2
+ 2\hat{\beta}\sum_{t=1}^T Y_t
- T\hat{\beta}^2 \\
&=
-2T\hat{\beta}\beta_0
+ T\beta_0^2
+ 2T\hat{\beta}^2
- T\hat{\beta}^2 \\
&=
T(\hat{\beta}-\beta_0)^2.
\end{align*}

Therefore,
\[
\lambda^*
=
(T-1)\frac{T(\hat{\beta}-\beta_0)^2}{\sum_{t=1}^T (Y_t-\hat{\beta})^2}.
\]

Now define the unbiased sample variance
\[
\hat{\sigma}_u^2
=
\frac{1}{T-1}\sum_{t=1}^T (Y_t-\hat{\beta})^2.
\]

Then
\[
\sum_{t=1}^T (Y_t-\hat{\beta})^2
=
(T-1)\hat{\sigma}_u^2,
\]
so that
\[
\lambda^*
=
\frac{T(\hat{\beta}-\beta_0)^2}{\hat{\sigma}_u^2}
=
\frac{(\hat{\beta}-\beta_0)^2}{\hat{\sigma}_u^2/T}.
\]

\paragraph{Distribution of the test statistic.}

Under the normality assumption, we have
\[
\frac{\sqrt{T}(\hat{\beta}-\beta)}{\sigma}
\sim N(0,1),
\qquad
\frac{(T-1)\hat{\sigma}_u^2}{\sigma^2}
\sim \chi^2(T-1),
\]
and these two quantities are independent.

Under $H_0$, it follows that
\[
\lambda^*
=
\frac{\left[\frac{\sqrt{T}(\hat{\beta}-\beta_0)}{\sigma}\right]^2}
{\frac{(T-1)\hat{\sigma}_u^2}{\sigma^2}/(T-1)}.
\]

Hence, $\lambda^*$ is the ratio of two independent chi-square variables,
so that
\[
\lambda^* \sim F(1,T-1).
\]

\paragraph{Equivalent $t$-statistic.}

Taking the square root,
\[
t
=
\frac{\hat{\beta}-\beta_0}{\sqrt{\hat{\sigma}_u^2/T}},
\]
we obtain
\[
t \sim t(T-1).
\]

\paragraph{Test procedure.}

\begin{enumerate}
\item Choose significance level $\alpha$.
\item Compute
\[
t = \frac{\hat{\beta}-\beta_0}{\sqrt{\hat{\sigma}_u^2/T}}.
\]
\item Find the critical value $t_{\alpha/2,T-1}$.
\item Reject $H_0$ if $|t| \ge t_{\alpha/2,T-1}$.
\end{enumerate}
\section{Example 2: Variance Test}

Let $(Y_1,\dots,Y_T)$ be a random sample from
\[
Y_t \sim N(\beta,\sigma^2).
\]

We want to test
\[
H_0:\sigma^2=\sigma_0^2,
\qquad
H_1:\sigma^2\neq\sigma_0^2,
\]
where $\sigma_0^2$ is known.

\paragraph{Likelihood function.}

The likelihood function of the sample is
\[
L(\beta,\sigma^2)
=
(2\pi\sigma^2)^{-T/2}
\exp\!\left(
-\frac{1}{2\sigma^2}
\sum_{t=1}^T (Y_t-\beta)^2
\right).
\]

The maximum likelihood estimators are
\[
\hat{\beta} = \frac{1}{T}\sum_{t=1}^T Y_t,
\qquad
\hat{\sigma}^2 = \frac{1}{T}\sum_{t=1}^T (Y_t-\hat{\beta})^2.
\]

\paragraph{Maximized likelihoods.}

The maximized likelihood functions are
\[
l^*(\Omega)
=
(2\pi\hat{\sigma}^2)^{-T/2} e^{-T/2},
\]
and under $H_0$,
\[
l^*(\Omega_0)
=
(2\pi\sigma_0^2)^{-T/2}
\exp\!\left(
-\frac{T}{2}\frac{\hat{\sigma}^2}{\sigma_0^2}
\right).
\]

\paragraph{Why these two expressions are different.}

In the unrestricted case, both $\beta$ and $\sigma^2$ are chosen freely to maximize the likelihood.
Thus, after substituting the maximum likelihood estimators
\[
\hat{\beta} = \frac{1}{T}\sum_{t=1}^T Y_t,
\qquad
\hat{\sigma}^2 = \frac{1}{T}\sum_{t=1}^T (Y_t-\hat{\beta})^2,
\]
we obtain
\[
\sum_{t=1}^T \frac{(Y_t-\hat{\beta})^2}{\hat{\sigma}^2}
=
\frac{\sum_{t=1}^T (Y_t-\hat{\beta})^2}
{\frac{1}{T}\sum_{t=1}^T (Y_t-\hat{\beta})^2}
=
T.
\]
Therefore,
\[
\exp\!\left(
-\frac{1}{2\hat{\sigma}^2}
\sum_{t=1}^T (Y_t-\hat{\beta})^2
\right)
=
\exp(-T/2),
\]
which explains the simplified form of $l^*(\Omega)$.

Under the null hypothesis, however, $\sigma^2$ is no longer free: it is fixed at the known value $\sigma_0^2$.
Only $\beta$ is estimated, and its restricted maximum likelihood estimator is still
\[
\hat{\beta} = \frac{1}{T}\sum_{t=1}^T Y_t.
\]
Hence the exponent becomes
\[
-\frac{1}{2\sigma_0^2}\sum_{t=1}^T (Y_t-\hat{\beta})^2.
\]
Using
\[
\sum_{t=1}^T (Y_t-\hat{\beta})^2 = T\hat{\sigma}^2,
\]
this can be written as
\[
-\frac{1}{2\sigma_0^2}\sum_{t=1}^T (Y_t-\hat{\beta})^2
=
-\frac{1}{2\sigma_0^2}\,T\hat{\sigma}^2
=
-\frac{T}{2}\frac{\hat{\sigma}^2}{\sigma_0^2}.
\]
Therefore,
\[
l^*(\Omega_0)
=
(2\pi\sigma_0^2)^{-T/2}
\exp\!\left(
-\frac{T}{2}\frac{\hat{\sigma}^2}{\sigma_0^2}
\right).
\]

The key point is that in the unrestricted likelihood, $\sigma^2$ is replaced by its own maximum likelihood estimator, so the exponent simplifies to $-T/2$.
Under $H_0$, by contrast, $\sigma^2$ is fixed at $\sigma_0^2$, so the ratio $\hat{\sigma}^2/\sigma_0^2$ remains in the exponent.


\paragraph{Likelihood ratio.}

The likelihood ratio is
\[
\lambda
=
\frac{l^*(\Omega)}{l^*(\Omega_0)}
=
\left(\frac{\hat{\sigma}^2}{\sigma_0^2}\right)^{-T/2}
\exp\!\left[
-\frac{T}{2}\left(1-\frac{\hat{\sigma}^2}{\sigma_0^2}\right)
\right].
\]

Since $\lambda$ is inconvenient to work with directly, define the statistic
\[
S
=
\frac{(T-1)\hat{\sigma}_u^2}{\sigma_0^2},
\]
where
\[
\hat{\sigma}_u^2
=
\frac{1}{T-1}\sum_{t=1}^T (Y_t-\hat{\beta})^2
\]
is the unbiased estimator of $\sigma^2$.

\paragraph{Distribution of the test statistic.}

Under the normality assumption, we have
\[
\frac{(T-1)\hat{\sigma}_u^2}{\sigma^2}
\sim \chi^2(T-1).
\]

Under $H_0:\sigma^2=\sigma_0^2$, it follows that
\[
S
=
\frac{(T-1)\hat{\sigma}_u^2}{\sigma_0^2}
\sim \chi^2(T-1).
\]

\paragraph{Interpretation.}

The likelihood ratio $\lambda$ is maximized when $S=T$, and becomes small when $S$ is either very large or very small.
Thus, the rejection region is two-sided.

\paragraph{Test procedure.}

\begin{enumerate}
\item Choose significance level $\alpha$.
\item Find critical values $S_L$ and $S_U$ such that
\[
P(S \le S_L)=\frac{\alpha}{2},
\qquad
P(S \ge S_U)=\frac{\alpha}{2},
\]
using the $\chi^2(T-1)$ distribution.
\item Compute
\[
S = \frac{(T-1)\hat{\sigma}_u^2}{\sigma_0^2}.
\]
\item Reject $H_0$ if $S \le S_L$ or $S \ge S_U$.
Otherwise, fail to reject $H_0$.
\end{enumerate}



\section{Example 3: Equality of Variances}

Let
\[
Y_1,\dots,Y_{T_1} \sim N(\beta_1,\sigma_1^2),
\qquad
X_1,\dots,X_{T_2} \sim N(\beta_2,\sigma_2^2).
\]

Test
\[
H_0:\sigma_1^2=\sigma_2^2,
\qquad
H_1:\sigma_1^2\neq\sigma_2^2.
\]

Define
\[
\hat{\sigma}_1^2
=
\frac{1}{T_1-1}\sum_{t=1}^{T_1}(Y_t-\bar{Y})^2,
\qquad
\hat{\sigma}_2^2
=
\frac{1}{T_2-1}\sum_{t=1}^{T_2}(X_t-\bar{X})^2.
\]

The test statistic is
\[
F
=
\frac{\hat{\sigma}_1^2}{\hat{\sigma}_2^2}.
\]

Under $H_0$,
\[
F \sim F(T_1-1,T_2-1).
\]

\paragraph{Test procedure.}
\begin{enumerate}
\item Choose $\alpha$.
\item Find critical values $F_L, F_U$.
\item Compute $F$.
\item Reject $H_0$ if $F \le F_L$ or $F \ge F_U$.
\end{enumerate}

\section{Asymptotic Tests}

\subsection{Likelihood Ratio (LR)}

\[
LR = 2\ln\left(\frac{l^*(\Omega)}{l^*(\Omega_0)}\right)
\]

\[
LR \xrightarrow{d} \chi^2(J)
\]

\subsection{Wald Test}

\[
W = R(\hat{\theta})'
\left[
\frac{\partial R}{\partial \theta}
V(\hat{\theta})
\frac{\partial R'}{\partial \theta}
\right]^{-1}
R(\hat{\theta})
\]

\[
W \xrightarrow{d} \chi^2(J)
\]

\subsection{Lagrange Multiplier (LM)}

\[
LM =
\left(\frac{\partial \ln l}{\partial \theta}\right)
I(\theta)^{-1}
\left(\frac{\partial \ln l}{\partial \theta}\right)'
\]

\[
LM \xrightarrow{d} \chi^2(J)
\]

\section{Normal Mean Example}

\[
H_0: \beta = k
\]

\subsection{LR}

\[
LR = T \ln\left(\frac{\hat{\sigma}_0^2}{\hat{\sigma}^2}\right)
\]

\subsection{Wald}

\[
W = \frac{(\hat{\beta} - k)^2}{\hat{\sigma}^2/T}
\]

\subsection{LM}

\[
LM = \frac{(\hat{\beta} - k)^2}{\hat{\sigma}_0^2/T}
\]

\section{Relation Between Tests}

\[
W \ge LR \ge LM
\]

Asymptotically:
\[
LR, W, LM \sim \chi^2(J)
\]